\documentclass[12pt]{article}
\usepackage[utf8]{inputenc}
\usepackage[T1]{fontenc}
\usepackage[margin=1in]{geometry}
\usepackage{setspace}
\usepackage{xcolor}
\usepackage[colorlinks=true, urlcolor=blue!60!black, linkcolor=blue!60!black]{hyperref}
\setlength{\parindent}{0em}
\setlength{\parskip}{0.6em}
\renewcommand{\baselinestretch}{1.0}
\begin{document}
\begin{center}
{\LARGE \textit{Ex parte Milligan}, 71 U.S. (4 Wall.) 2 (1866)\par}
\vspace{0.5em}
{\large Case Brief --- Two-Page Edition\par}
\end{center}

\section{Memory Jogger}
An Indiana civilian sentenced to hang by a Civil War military commission won release: unanimously, the commission lacked jurisdiction; 5--4, neither the President nor Congress may try civilians by military commission where civil courts are open, even with habeas suspended (\href{https://supreme.justia.com/cases/federal/us/71/2/}{71 U.S. (4 Wall.) 2}).

\section{Detailed Case Facts}
Congress's 1863 Habeas Corpus Act allowed suspension of the writ but required discharge of civilian detainees in loyal States if a federal grand jury adjourned without indicting them. Lambdin Milligan, a lifelong civilian in Indiana, was arrested in 1864 and sentenced to hang by a military commission for aiding the ``Sons of Liberty'' plot, while the federal court in Indianapolis sat and its grand jury did not indict him.

\section{Procedural History}
Milligan sought habeas corpus in the Circuit Court, whose divided judges certified three questions to the Supreme Court.

\section{Judicial Votes}
9--0 for Milligan; 5--4 on Congress's power. Davis wrote for Nelson, Grier, Clifford, and Field; Chase, with Wayne, Swayne, and Miller, concurred in the judgment.

\section{Holding}
The commission lacked jurisdiction. Neither the President nor Congress may subject civilians unconnected to the armed forces to military trial in States where the civil courts are open and unobstructed.

\section{Analysis of Opinions}
Davis held that the Constitution governs in war as in peace and that only the privilege of habeas corpus may be suspended. Article III and the Fifth and Sixth Amendments guarantee civilians indictment and jury trial, and martial rule is confined to places where war has actually closed the courts. Suspension permits detention, but ``[t]he Constitution goes no further.'' Chase agreed that the 1863 Act required release but argued that Congress could authorize commissions where courts, though open, are ``wholly incompetent.'' \textit{Ex parte Quirin} later allowed commission trials of unlawful belligerents, even citizens.

\section{Examples: Future Applications}
\textbf{Same-Side:} a President's martial-law tribunal for a student arsonist while courts are open is invalid. \textbf{Opposite-Side:} a citizen who joins enemy forces and commits sabotage may be tried by commission under \textit{Quirin}. \textbf{Fence-Sitter:} commissions authorized by Congress where insurrectionists' threats keep an open court from seating juries.

\section{Critique}
The Court acted only after the war, having stayed its hand in \textit{Vallandigham}, and Chase showed that courts may be open yet compromised. \textit{Quirin}'s belligerency exception leaves a hard boundary between protected civilians and triable saboteurs.

\section{Key Quotations}
\begin{quote}
``The Constitution of the United States is a law for rulers and people, equally in war and in peace, and covers with the shield of its protection all classes of men, at all times and under all circumstances. No doctrine involving more pernicious consequences was ever invented by the wit of man than that any of its provisions can be suspended during any of the great exigencies of government.''
\end{quote}
\noindent --- Justice Davis, Opinion of the Court, 71 U.S. at 120--21.

\end{document}
